% ==============================================================================
% CAPÍTULO 1 - INTRODUÇÃO
% ==============================================================================

\chapter{Introdução}
\label{cap:introducao}

A transformação digital tem redefinido a gestão operacional em diversos setores de prestação de serviços técnicos, desde a manutenção eletrônica até o reparo automotivo. Nesse contexto, a eficiência na gestão de ordens de serviço, o controle rigoroso de estoque e a conformidade tributária tornaram-se pilares indispensáveis para a sustentabilidade de pequenos e médios estabelecimentos. A FK Oficina, consolidada como uma provedora de soluções \textit{Software as a Service} (SaaS), atende a essa demanda crescente por meio de um ecossistema digital integrado.

Contudo, a expansão de uma base de clientes heterogênea — abrangendo tanto oficinas mecânicas quanto assistências técnicas de eletrônicos — impôs desafios complexos de infraestrutura. A necessidade de gerenciar centenas de instâncias isoladas de bancos de dados (\textit{multi-tenancy}), automatizar o licenciamento financeiro e sincronizar tabelas tributárias nacionais exigiu a criação de uma camada de gestão superior.

Este relatório documenta o desenvolvimento do projeto \textbf{Oficina Master}, um \textit{hub} administrativo centralizado projetado para orquestrar o ecossistema da FK Oficina. O projeto visou mitigar a carga operacional da equipe técnica e unificar o controle sobre os \textit{tenants}, permitindo uma visão holística e automatizada do ecossistema.

%-------------------------------------------------
% Justificativa
%-------------------------------------------------
\section{Justificativa}
\label{sec:justificativa}

A relevância deste trabalho fundamenta-se na otimização da eficiência operacional técnica. Antes da implementação do Oficina Master, rotinas críticas como migrações de esquemas de banco de dados e atualizações do Índice Brasileiro de Planejamento e Tributação (IBPT) eram realizadas manualmente, resultando em alto custo temporal e risco de inconsistências. Academicamente, o estágio proporcionou a aplicação prática de padrões de projeto avançados, como Injeção de Dependências e Arquitetura Orientada a Serviços (SOA), consolidando as competências previstas no currículo de Ciência da Computação.

%-------------------------------------------------
% Objetivos
%-------------------------------------------------
\section{Objetivos}
\label{sec:objetivos}

\subsection{Objetivo Geral}
\label{subsec:objective-geral}

Desenvolver e implementar um \textit{hub} administrativo centralizado (\textit{Oficina Master}) para a orquestração de um ecossistema SaaS, automatizando processos de gestão de dados, licenciamento e suporte técnico.

\subsection{Objetivos Específicos}
\label{subsec:objetivos-especificos}

\begin{alineas}
    \item Projetar e desenvolver uma API RESTful utilizando o \textit{framework} Laravel 12;
    \item Implementar uma interface reativa utilizando Vue.js 3 com gerenciamento de estado via Pinia;
    \item Desenvolver um mecanismo de alternância dinâmica de contexto para operações em massa sobre múltiplos bancos de dados;
    \item Integrar o fluxo de licenciamento com a API do Mercado Pago, garantindo a automação do ciclo de vida das contas;
    \item Implementar rotinas de sincronização assíncrona para tabelas tributárias nacionais (IBPT);
    \item Estabelecer um painel de telemetria (\textit{dashboard}) para monitoramento proativo da integridade do ecossistema.
\end{alineas}

%-------------------------------------------------
% Estrutura do Trabalho
%-------------------------------------------------
\section{Estrutura do Trabalho}
\label{sec:estrutura}

Este relatório está organizado em seis capítulos. Além desta Introdução, o \autoref{ch:apresentacao-empresa} contextualiza a FK Oficina e sua trajetória. O \autoref{cap:metodologia} descreve os métodos ágeis e ferramentas de desenvolvimento. O \autoref{cap:revisao-bibliografica} fundamenta tecnicamente as tecnologias empregadas. O \autoref{cap:desenvolvimento} detalha a implementação e os desafios superados. Finalmente, o \autoref{cap:conclusao} sintetiza as considerações finais e os aprendizados obtidos.
